\section{Automated post-training as continual domain adaptation}

The human-learning analogy suggests a closed loop rather than one final training run. A
learner studies available material, tests themselves, identifies what remains unknown,
builds targeted practice, and tests again. For an LLM agent, the corresponding loop begins
with an externally anchored objective or domain. The system creates diagnostic tasks,
attempts them, analyzes failures, synthesizes a curriculum, and repeats. Post-training
becomes continual learning when verified experience changes the policy while earlier
capabilities are retained, rather than merely appending another isolated fine-tune.

Figure~\ref{fig:auto-post-training-loop} captures the proposed control loop. Its most
important decision is the reflection step: a failed attempt may reveal a local mistake, a
missing capability, a harness defect, or an unreliable evaluator. Only a genuine capability
gap should trigger training. The corrected loop splits curriculum tasks from protected
evaluation, converts diagnosed gaps into verified demonstrations, rewards, or textual
hints, updates a versioned policy through SFT, RL, or distillation, and promotes the
candidate only after new-domain and retention tests pass.

\begin{figure}[p]
  \centering
  \resizebox{0.96\textwidth}{!}{%
  \begin{tikzpicture}[
      node distance=4mm and 8mm,
      flow/.style={-Latex, semithick, draw=black!70},
      loop/.style={-Latex, semithick, draw=blue!60!black},
      oversight/.style={-Latex, semithick, dashed, draw=violet!70},
      box/.style={draw=black!60, rounded corners=2pt, align=center,
                  font=\scriptsize, minimum height=7.5mm, text width=4.2cm,
                  inner sep=2.5pt},
      anchorbox/.style={box, fill=violet!10, draw=violet!65, text width=5.4cm},
      taskbox/.style={box, fill=cyan!8, draw=cyan!55!black, text width=3.9cm},
      diagnosebox/.style={box, fill=yellow!14, draw=yellow!45!black},
      trainbox/.style={box, fill=green!9, draw=green!50!black},
      evalbox/.style={box, fill=blue!7, draw=blue!55!black},
      sidebox/.style={box, fill=black!3, text width=3.3cm},
      label/.style={font=\scriptsize, fill=white, inner sep=1pt}
    ]
    \node[anchorbox] (objective)
      {Externally anchored objective, domain sources, and success criteria};
    \node[box, below=of objective] (current) {Current policy $\pi_k$ + agent harness};
    \node[taskbox, below=of current, text width=5.4cm] (generate)
      {Generate diagnostic tasks from observed failures and target capabilities};

    \node[taskbox, below left=6mm and 4mm of generate] (curriculum)
      {Adaptation pool\\training and practice tasks};
    \node[evalbox, below right=6mm and 4mm of generate, text width=3.9cm] (protected)
      {Protected evidence\\domain tests + retention suite};

    \node[box, below=7mm of curriculum] (attempt)
      {Attempt adaptation tasks and record exact traces};
    \node[diagnosebox, below=of attempt] (diagnose)
      {Verify outcomes and diagnose root cause};
    \node[sidebox, left=7mm of diagnose] (local)
      {Local mistake, harness defect, or evaluator failure\\repair and retry without training};
    \node[trainbox, below=of diagnose] (gap)
      {Evidence of a genuine capability or knowledge gap};
    \node[trainbox, below=of gap] (signals)
      {Synthesize and validate supervision\\demonstrations $\mid$ rewards $\mid$ semantic hints};
    \node[trainbox, below=of signals] (posttrain)
      {Continual post-training\\SFT $\mid$ RL $\mid$ OPD $\mid$ MOPD};
    \node[box, below=of posttrain] (candidate) {Candidate policy $\pi_{k+1}$};

    \node[evalbox, right=9mm of candidate] (evaluate)
      {Evaluate candidate on protected domain and retention tests};
    \node[evalbox, below=of evaluate] (decision)
      {Verified improvement with acceptable retention?};
    \node[sidebox, below left=6mm and 5mm of decision] (rollback)
      {No\\rollback and revise diagnosis, data, or objective};
    \node[trainbox, below right=6mm and 5mm of decision, text width=3.3cm] (promote)
      {Yes\\promote $\pi_{k+1}$ and begin the next gap-discovery cycle};

    \node[anchorbox, above=5mm of objective] (human)
      {Human oversight at objective, uncertain diagnoses, and promotion gates};

    \draw[flow] (objective) -- (current);
    \draw[flow] (current) -- (generate);
    \draw[flow] (generate) -- (curriculum);
    \draw[flow] (generate) -- (protected);
    \draw[flow] (curriculum) -- (attempt);
    \draw[flow] (attempt) -- (diagnose);
    \draw[flow] (diagnose) -- node[label, right] {capability gap} (gap);
    \draw[flow] (gap) -- (signals);
    \draw[flow] (signals) -- (posttrain);
    \draw[flow] (posttrain) -- (candidate);
    \draw[flow] (candidate) -- (evaluate);
    \draw[flow] (protected) |- (evaluate);
    \draw[flow] (evaluate) -- (decision);
    \draw[flow] (decision) -- (rollback);
    \draw[flow] (decision) -- (promote);

    \draw[flow] (diagnose) --
      node[label, midway, above, align=center, text width=1.35cm] {not a\\model gap}
      (local);
    \draw[loop] (local.north) |- node[label, near start] {repair} (attempt.west);
    \draw[loop] (rollback.west) -| (diagnose.south west);
    \draw[loop] (promote.east) |- (current.east);

    \draw[oversight] (human) -- (objective);
    \draw[oversight] (human.east) -- ++(4mm,0) |- (diagnose.east);
    \draw[oversight] (human.east) -- ++(4mm,0) |- (decision.east);

    \node[draw=black!35, dashed, rounded corners=2pt,
          fit=(curriculum)(protected), inner sep=3mm] (split) {};
    \node[label, anchor=south] at (split.north)
      {adaptation data and evaluation evidence remain separate};
  \end{tikzpicture}
  }
  \caption{Proposed automated continual post-training loop. Diagnosed capability gaps are
  converted into validated supervision and a versioned policy update. The candidate is
  promoted only if it improves a protected domain evaluation while retaining earlier
  capabilities; otherwise the system rolls back and revises the diagnosis or curriculum.
  This is a conceptual proposal, not a reproduced training system.}
  \label{fig:auto-post-training-loop}
\end{figure}


Self-benchmarking creates the main reliability problem. If one model generates the tasks,
answers them, judges them, and trains on them, the loop can reward its own misconceptions,
leak evaluation items into training, or construct an increasingly easy curriculum. A
credible loop therefore separates training tasks from held-out tests, grounds verification
in an environment or independent judge, records data provenance, evaluates regression on
earlier capabilities, and retains a human gate or rollback path when evidence is uncertain.
Self-distillation has shown that sequential skills can be acquired with less forgetting
than ordinary SFT, but this is evidence for one update mechanism, not yet for a fully
autonomous improvement loop \citep{shenfeld2026selfdistillation}.

This distinction also separates two ambitions. Automated domain adaptation is relatively
well anchored: a human supplies the target domain, source material, task constraints, and
often verifiable tests. Recursive self-improvement is stronger because the system must also
improve its gap detector, curriculum builder, and verifier without losing contact with an
external objective. Progress should therefore mean improvement on protected, independently
measured tasks with capability retention---not merely higher reward on the system's own
generated benchmark.

The earlier methods occupy different roles inside this loop. SFT clones verified examples;
RL searches against outcome evidence; critics and semantic feedback locate errors; OPD
transfers a correction on student-visited states; and MOPD merges specialized teachers.
None closes the loop alone. The position of this paper is that long-term agent improvement
requires their orchestration around trustworthy gap detection, hierarchical feedback,
continual evaluation, and retention.
